\documentclass[12pt,a4paper]{extarticle} % Используем расширенный класс

\usepackage[T2A]{fontenc}
\usepackage[utf8]{inputenc}
\usepackage[english, russian]{babel}

% Улучшенный Times-образный шрифт
\usepackage{tempora}
\usepackage{newtxmath}
\usepackage{microtype} % Чтобы текст ложился ровно и красиво
\usepackage{type1ec}

% Настройка полей строго по Приложению №4 (к городской НПК)
\usepackage[
left=30mm,
right=10mm,
top=20mm,
bottom=20mm,
footskip=10mm
]{geometry}

\usepackage[nodisplayskipstretch]{setspace}
\onehalfspacing 

\usepackage{indentfirst}
\setlength{\parindent}{1.25cm}

% Настройка заголовков: центрирование, жирный
\usepackage{titlesec}
\titleformat{\section}{\normalfont\Large\bfseries\centering}{}{0pt}{}
\titleformat{\subsection}{\normalfont\large\bfseries\centering}{}{0pt}{}

\usepackage{graphicx}
\usepackage{caption}

% Нумерация страниц в центре нижнего поля
\usepackage{fancyhdr}
\pagestyle{fancy}
\fancyhf{}
\fancyfoot[C]{\thepage}
\renewcommand{\headrulewidth}{0pt}

%Правильные сниппеты с кодом
\usepackage{minted}
\setminted{
	breaklines=true,
	breakanywhere=true,
	fontsize=\footnotesize,
	baselinestretch=1.0,  % Устанавливает одинарный интервал внутри кода
	breaksymbol={},           % Удаляем символ переноса
	fontfamily=tt,
	xleftmargin=-2.15cm,  % Делает отступ как у красной строки
}
% Отступы у списков
\usepackage{enumitem}

% Настройка для всех маркированных списков
\setlist[itemize]{
	noitemsep,           % Убирает интервал между пунктами
	topsep=0pt,          % Убирает отступ перед списком
	parsep=0pt,          % Убирает отступ внутри пункта (если текст в несколько строк)
	partopsep=0pt,
	leftmargin=1.25cm    % Выравнивание по красной строке
}

% Настройка для нумерованных списков
\setlist[enumerate]{
	noitemsep,
	topsep=4pt,
	leftmargin=1.25cm
}
% Настройка гиперссылок в документе
\usepackage[unicode, pdftex]{hyperref}
\usepackage{url}

\hypersetup{
	citecolor=green,       % Цвет ссылок на литературу
	filecolor=blue,     
	urlcolor=blue,         % Цвет внешних URL-ссылок
	pdftitle={Разработка веб-сайта для нахождения оптимальной рассадки учеников в классе}, % Метаданные PDF
	pdfauthor={Прокофьев Даниил}
}

% Картинки
\usepackage{graphicx}
\usepackage{float}
\graphicspath{{../presentation/public/screenshots}}

 % =================================================================
\begin{document}
	
	% --- ТИТУЛЬНЫЙ ЛИСТ ---
	\begin{titlepage}
		\begin{center}
			\setstretch{1.0}
			НПК «Интерес. Познание. Творчество» \\
			\vspace{4cm}
			
			Секция: Информатика \\
			\vspace{1cm}
			
			{\Large \textbf{Разработка веб-сайта для нахождения оптимальной рассадки учеников в классе}} \\
			\vfill
			
			\begin{flushright}
				\begin{minipage}{0.5\textwidth}
					\textbf{Автор работы:} \\
					Прокофьев Даниил \\
					МБОУ Гимназия, 10 класс \\
					\vspace{0.5cm}
					
					\textbf{Руководитель:} \\
					Елена Васильевна Утянская, \\
					учитель информатики
				\end{minipage}
			\end{flushright}
			
			\vfill
			г. Обнинск \\
			2026
		\end{center}
	\end{titlepage}
	
	\newpage
	\setcounter{page}{2}
	
	% --- ОГЛАВЛЕНИЕ ---
	\tableofcontents
	\newpage
	
	% --- ВВЕДЕНИЕ ---
	\section*{Введение}
	\addcontentsline{toc}{section}{Введение}
	У подавляющего большинства классных руководителей рассадку своего класса принято делать вручную. Но эффективно ли это? Учитывает ли такая ручная рассадка предпочтения всех и каждого?
	
	Вряд ли. Американским психологом (Miller 1956) Джорджем Миллером было выяснено, что человек не способен держать в своей кратковременной памяти более $7 \pm 2$  факторов одновременно \cite{miller1956}.
	
	Это значит, что оптимальность рассадки, создаваемой учителем, снижается с увеличением количества учеников и требований к рассадке. Если рассадить до 12 человек с не более чем 5-7 ограничениями оптимально для учителя не составляет никакого труда, то класс на 30 человек с его десятками перекрестными ограничениями и предпочтениями будет рассажен плохо.
	
	Целью данной работы является изучение способов решения задачи рассадки и учеников и написание веб-сайта, который поможет учителям быстро и эффективно рассаживать учащихся.
	
	\newpage
	
	% --- ОСНОВАНАЯ ЧАСТЬ ---
	\section{Выбор технологий и архитектуры}
	
	Для реализации системы была выбрана клиент-серверная архитектура, разделяющая интерфейс пользователя (\textit{frontend}) и вычислительный модуль (\textit{backend}). 
	
	Вынос ресурсоемких вычислений на сторону сервера продиктован двумя факторами:
	\begin{itemize}
		\item \textbf{Аппаратные ограничения:} Недостаточная мощность компьютеров учителей в образовательных учреждениях
		\item \textbf{Производительность:} Согласно исследованиям \cite{usenix_webassembly}, нативные вычисления значительно превосходят интерпретируемые аналоги в браузере
	\end{itemize}
	
	\subsection{Вычисления}
	Выбор языка \texttt{Go} для реализации алгоритма рассадки обусловлен следующими характеристиками:
	\begin{itemize}
		\item \textbf{Эффективность:} Как компилируемый язык, \texttt{Go} показывает производительность, кратно превосходящую \texttt{Python} \cite{benchmarks_game}
		\item \textbf{Параллелизм:} Использование механизмов \textit{goroutines} позволяет эффективно распараллеливать вычисления
		\item \textbf{Надежность:} Строгая статическая типизация минимизирует вероятность ошибок при работе со сложными структурами данных.
		\item \textbf{Портативность:} Компиляция в единый бинарный файл упрощает развертывание системы
	\end{itemize}
	
	\subsection{Пользовательский интерфейс}
	В качестве фронтенд-фреймворка используется \texttt{Vue.js}. Несмотря на то, что вопрос выбора фреймворка максимально субъективен, так как каждый выбирает то, что ему нравится и идеального решения не существует, он обладает такими плюсами, как реактивность и компонентный подход, что позволяет облегчить разработку приложения.
	
	% --- ОСНОВНАЯ ЧАСТЬ ---
	\section{Формализация задачи алгоритма}
	\label{formalization}
	Чтобы определить, является ли решение правильным, нужно придумать критерии его оценивания.
	
	Было выделено четыре типа ограничений и предпочтений, которые удовлетворят желания большинства пользователей и не ограничивают его свободу выбора:
	
	\begin{enumerate}
		\item \textbf{Медицинские предписания} (например, ученик должен сидеть за первой партой согласно предписанию врача)
		\item \textbf{Нежелательное соседство} (например, два ученика враждуют)
		\item \textbf{Предпочтения по рядам и партам} (например, ученик хочет сидеть на первой парте)
		\item \textbf{Предпочтения учеников по соседству} (например, два друга хотят сидеть вместе)
	\end{enumerate}
	
	\section {Выбор алгоритма}
	
	Задача относится к классу \textbf{NP-полных}. Это значит, что не существует алгоритма, который за разумное (полиноминальное) время сумел бы решить задачу. Тем не менее, есть эвристические алгоритмы, которые способны найти приближенный ответ, иногда близкий к наилучшему возможному. Рассмотрим некоторые из тех, что применимы в нашем случае.
	
	\subsection{Жадный алгоритм}
	Для решения задачи можно попробовать использовать алгоритм, который наиболее удобен при ручной рассадке учеников. В чем его суть?
	
	\begin{itemize}
		\item Выставим учеников в очередь. Сначала поставим учеников, имеющих медицинские ограничения, потом учеников, имеющих предпочтения по ряду или парте, потом тех, кому вместе сидеть нельзя и, наконец друзей (порядок предпочтений можно менять)
		\item  Возьмем одного ученика из очереди и посадим его на лучшее в данный момент для него место (принцип жадности)
	\end{itemize}
	
	Однако, например, посаженный на первых порах неудачно, ученик будет мешать получению оптимальной рассадки. Алгоритм найдет решение, но это будет локальный максимум, а не глобальный, который мы ищем.
	
	\subsection{Метод отжига}
	Отжиг в физике — это процесс сильного нагревания металла, за которым следует его медленное охлаждение. Идея состоит в том, чтобы выстроить атомы в кристаллическую решётку и сформировать как можно больше таких устойчивых решёток в металле. Если всё сделать правильно, у металла снижается твёрдость, убираются внутренние напряжения и металл становится более однородным — так его легче обрабатывать.
	
	Если выписать из процесса отжига ключевые мысли и точки, то получим идею к реализации алгоритма \cite{otzhig-kod}:
	\begin{itemize}
		\item есть система, которая находится не в нужном для нас состоянии
		\item есть процессы, которые могут протекать в этой системе
		\item есть параметр (температура), благодаря которому система сама может отрегулировать себя и прийти в нужное состояние
		\item и есть конечное значение (температуры), по достижении которого система считается максимально близкой к нужному состоянию
	\end{itemize}
	
	Однако сам алгоритм довольно медленный, так как для получения оптимального результата требует медленного охлаждения. А при быстром охлаждении алгоритм все так же может застревать в локальном максимуме.
	
	\subsection{Генетический алгоритм}
	Обратимся к другому алгоритму - генетическому. Будем развивать не одно решение, а целую популяцию, и будем производить над ней такие операции как естественный отбор, скрещивание, случайная мутация.
	
	Однако в данном случае эволюция получается слепая. То есть мутации часты, но случайны, и на то, чтобы переместить двух учеников местами, алгоритму может потребоваться очень много поколений, что увеличивает вычислительную сложность.
	
	\subsection{Меметический алгоритм}
	Генетический алгоритм можно усовершенствовать. Добавим к эволюции механизм “обучения” особи.
	
	Будем к наилучшей особи и к случайной с шансом 5\% применять локальный поиск. Локальный поиск будет совершать некое число итераций, пытаясь на каждой из них поменять двух случайных учеников так, чтобы качество рассадки улучшилось. Если после перемещения рассадка улучшилась, ты мы оставляем новую версию, если нет, то возвращаем все, как было.
	
	Этот подход предотвращает застревание в локальном максимуме и помогает в особо тяжелых случаях.
	
	\section{Написание меметического алгоритма}
	
	\subsection{Как хранить рассадки в памяти?}
	Были использованы одномерные массивы, поскольку с ними любые операции производить гораздо легче. Однако возникает проблема - если учеников в классе меньше чем мест, то для того, чтобы ученик мог занять любое место в классе необходимо, чтобы одномерный массив был размера большего, чем количество учеников.
	
	Каждому ученику присваивается номер от 1 до n (количества учеников). Все, что больше n, считается как пустое место (-1).
	
	\begin{minted}{text}
		[1, 3, 5, 6, 7, 2, 4, 8, 9, 10] — класс размером 2 ряда * 5 парт в ряду, но только 8 учеников.
	\end{minted}
	
	Можно представить как:
	
	\begin{minted}{text}
		1        3
		5        6
		7        2
		4        8
		9(-1)   10(-1)
	\end{minted}
	
	\subsection{Как работает меметический алгоритм?}
	Рассмотрим сам цикл популяций:
	
	\begin{enumerate}
		\item Идем по поколениям, число которых высчитывается по формуле:
		\begin{equation}
			Generation = \min(300 + 20 \cdot n, 2500),
		\end{equation}
		где $n$ --- количество учеников.
		\item Оценим каждую особь с помощью функции \texttt{fitness()}
		\item Найдем индекс лучшей рассадки в \texttt{scores}
		\item Проверяем для предотвращения стагнации - если за последние поколение результат не улучшился более чем на 0.0001, прибавляем 1 к \texttt{stagnationCounter}
		\item Если \texttt{stagnationCounter >= stagnationLimit} (максимального количества поколений, в течение которых позволено оценке лучшей рассадки не меняться достаточно сильно в лучшую сторону), то прерываем программу
		\item Встряхиваем популяцию, если алгоритм близок к стагнации
		\item Копируем лучшую особь в новую популяцию и применяем к ней локальный поиск (элитизм)
		\item С помощью турнира отбираем две случайные особи и их скрещиваем
		\item С небольшим шансом в получившемся потомке меняем двух учеников случайным образом
	\end{enumerate}
	
	Код функции, запускающий алгоритм решения без вспомогательных подпрограмм представлен в Приложении 1.
	
	\subsection{Как реализовано проверка предпочтений?}
	Для проверки предпочтений были написаны вспомогательные функции: \texttt{checkMed}, \texttt{checkPref}, \texttt{checkFriends} и \texttt{checkEnemies}. Все они возвращают значение от 0.0 до 1.0 - насколько выполнено каждое из предпочтений ученика.
	
	Рассмотрим одну из них:
	
	\begin{minted}{go}
		func checkPref(student optStudent, row, col int) float64 {
			score := 0.0
			maxPossible := 0.0
			if len(student.pCols) > 0 {
				maxPossible += 1.0
			}
			if len(student.pRows) > 0 {
				maxPossible += 1.0
			}
			if maxPossible == 0 {
				return 0.0
			}
			if student.pCols[col] {
				score += 1.0
			}
			if student.pRows[row] {
				score += 1.0
			}
			return score / maxPossible
		}
	\end{minted}
	
	Видно, что функция считает сначала максимальное количество возможных баллов за выбранные ряд и парту (так как не у каждого могут быть такого вида предпочтения) и фактическое количество баллов за место. Затем делит одно на другое и получает число от 0 до 1. Это и будет учтённость предпочтения.
	
	Стоит сказать, что для ускорения работы программы был использован прием создания статического массива. Взяв предпочтения, которые не зависят от расположения учеников относительно друг друга, мы для каждого из них построили массив, который для каждого возможного места будет содержать количество баллов за место (учитывая предпочтения по месту и медицинские ограничения). Это намного быстрее, чем каждое поколение для каждого ученика проходить по его массиву парт или рядов.
	
	Для проверки друзей и нежелательного соседства используется матрица смежности. Они содержит записи вида:
	
	\begin{minted}{go}
		friendsMap[id1][id2] = true // если они хотят быть друзьями
	\end{minted}
	
	где id1 и id2 - номера учеников.
	
	\subsection{Веса фитнес-функции}
	Каждое значение функции проверки предпочтений домножается на число --- вес данного предпочтения. Это число, которое получено от пользователя, лежит в промежутке от 0.0 до 1.0 и показывает, насколько одно из предпочтений/ограничений важнее другого. Их 5: те, что перечислены в \ref{formalization} и параметр \texttt{Fill}, отвечающий за то, насколько ученики будут стремиться сесть ближе к первой парте
	
	Вот как это выглядит в запросе к серверу:
	\begin{minted}{json}
		...
		"PriorityWeights": {
			"Medical": 0.9,
			"Friends": 0.65,
			"Enemies": 0.7,
			"Preferences": 0.6,
			"Fill": 0.3
		}
		...
	\end{minted}
	
	\subsection{Функция fitness: как оценить рассадку?}
	Функция \texttt{fitness} считает оценку рассадки по такой формуле:
	\begin{equation}
		S_{total} = \frac{1}{N} \sum_{i=1}^{N} \frac{F_i \cdot W_{friend} - E_i \cdot W_{enemy} + S_{static, i, j}}{W_{friend} + W_{enemy} + W_{medical} + W_{pref}},
	\end{equation}
	
	где 
	\begin{itemize}
		\item $S_{total}$ --- общая приспособленность (fitness) варианта рассадки
		\item $N$ --- общее количество учеников в классе
		\item $F_i$ --- удовлетворенность $i$-го ученика по количеству друзей в радиусе близости
		\item $E_i$ --- неудовлетворенность $i$-го ученика по количеству врагов в радиусе близости
		\item $W_{friend}, W_{enemy}, W_{medical}, W_{pref}$ --- весовые коэффициенты значимости друзей и врагов, медицины и предпочитаемых рядов и парт
		\item $S_{static, i, j}$ --- сумма весов за статичные предпочтения (медицинские ограничения и предпочитаемые ряды и парты)
	\end{itemize}
	
	Код функции:
	
	\begin{minted}{go}
		func fitness(seating []int, config ClassConfig, w Weights, friends SocialMap, enemies SocialMap, staticScores []float64, nStudents int, friendsCount []int, enemiesCount []int) float64 {
			score := 0.0
			sumWeights := w.FriendBonus + w.EnemyPenalty + w.MedPenalty + w.PrefBonus + w.RowBonus
			if sumWeights == 0 {
				sumWeights = 1
			}
			for i, studentIdx := range seating {
				if studentIdx < 0 || studentIdx >= nStudents {
					continue
				}
				row, col := i/config.Columns, i%config.Columns
				fScore := checkFriends(studentIdx, seating, row, col, config, friends, nStudents, friendsCount)
				ePenalty := checkEnemies(studentIdx, seating, row, col, config, enemies, nStudents, enemiesCount)
				sScore := (fScore * w.FriendBonus) - (ePenalty * w.EnemyPenalty)
				sScore += staticScores[studentIdx*config.Rows*config.Columns+i]
				sScore /= sumWeights
				score += sScore
			}
			return score / float64(nStudents)
		}
	\end{minted}
	
	В данном случае функция показывает, насколько близка рассадка по качеству к идеальной --- такой, где у каждого ученика есть все виды предпочтений и они выполнены. То есть шкала оценки нормализована относительно теоретического максимума.
	
	Однако у большинства учеников набор активных ограничений и предпочтений ограничен, поэтому оценка рассадки может быть довольно небольшой --- в ходе тестирования минимальным значением было 20\%. Это не значит, что алгоритм плохо справился, а является эффектом выбора такой системы отсчета.
	
	\section{Пользовательский интерфейс}
	Стояла задача разработать понятный и удобный интерфейс. Нужно было отделить возможность сохранять классы, выбирать их в генераторе и изменять их параметры друг от друга, а также экспортировать файл рассадки в PDF, сохранять получившиеся рассадки и классы в памяти.
	Для этого был использован фреймворк \texttt{Vue.js} и библиотека \texttt{bootstrap-vue-next} (специальная версия Bootstrap, написанная для лучшей интеграции в Vue.js) и множество других пакетов, список которых можно посмотреть в приложении 2, содержащем файл \texttt{package.json} (в проектах на \texttt{Node.js} файл, где зафиксированы все пакеты, необходимые для сборки проекта)
	
	\begin{figure}[H]               % [H] - здесь"
		\centering                   % Выравнивание по центру
		\includegraphics[width=0.8\textwidth]{main.png} % 80% от ширины текста
		\caption{Главный экран программы} 
		\label{fig:main}        % Метка для ссылки в тексте
	\end{figure}
	
	\begin{figure}[H]
		\centering
		\includegraphics[width=0.8\textwidth]{editor.png}
		\caption{Интерфейс редактора класса} 
		\label{fig:editor}
	\end{figure}
	
	\begin{figure}[H]
		\centering
		\includegraphics[width=0.8\textwidth]{generator.png}
		\caption{Интерфейс страницы генерации} 
		\label{fig:generator}
	\end{figure}
	
	\section{Автоматизация развертывания и CI/CD}
	
	Для обеспечения стабильной работы и упрощения процесса деплоя приложения используется технология контейнеризации \textbf{Docker}. Она позволяет инкапсулировать серверную часть (как \textit{backend}, так и \textit{frontend}) со всеми необходимыми зависимостями, системными библиотеками и конфигурациями в единый изолированный образ.
	
	Применение Docker решает несколько критических задач:
	\begin{itemize}
		\item \textbf{Изоляция среды:} исключается влияние системных настроек сервера
		\item \textbf{Воспроизводимость:} гарантируется идентичность работы как на машине разработчика, так и на \textit{production}-сервере
		\item \textbf{Масштабируемость:} контейнерная архитектура позволяет быстро перезапускать или обновлять сервис без простоя системы.
	\end{itemize}
	
	Процесс сборки и доставки кода автоматизирован с помощью инструментов \textbf{GitHub Actions}. При каждом обновлении исходного кода в репозитории автоматически запускается сценарий (\textit{workflow}), выполняющий сборку образа контейнера и немного различающийся для \textit{frontend} и \textit{backend}.
	
	Такой подход минимизирует риск человеческой ошибки при обновлении и обеспечивает актуальность версии генератора рассадки на конечном сервере.
	
	\section{Веб-сайт в Интернете}
	Благодаря автоматической сборке несложно было развернуть свой собственный экземпляр генератора для демонстрации его возможностей, что я и сделал.
	
	Используя личный VPS сервер одного из крупных российских провайдеров и приобретенный домен \texttt{seating-generator.ru} был развернут экземпляр генератора, доступный для использования всем желающим.
	
	По сути, для развертывания всего лишь надо:
	
	\begin{itemize}
		\item Скачать \texttt{docker-compose.yml} файл из мета-репозитория \cite{github_repo} и, находясь в одной папке с этим файлом запустить команду \texttt{docker compose up --pull always}, которая скачает недостающие образы контейнеров и запустит их
		\item Настроить реверс-прокси, например \texttt{Caddy} или \texttt{Nginx} --- промежуточный сервер, принимающий запросы от пользователя и перенаправляющий их к генератору. Это необходимо прежде всего для безопасности, так как именно реверс-прокси берет на себя управление SSL-сертификатами, необходимыми для шифрования всех данных, пересылаемых между пользователем и сайтом и сокрытия "внутренней кухни" сервера
	\end{itemize}
	
	\section{Лицензирование}
	Проект полностью лицензирован под \textbf{GNU Affero Public License версии 3} \cite{license_agplv3}. Это значит:
	
	\begin{itemize}
		\item любой пользователь всегда будет иметь свободу изменять, распространять, неограниченно устанавливать, развертывать, использовать и изучать программу
		\item любое изменение к программе должно быть выпущено под той же самой лицензией, что гарантирует свободность программы
		\item если сервис доступен через Интернет, исходный код обязан быть открытым (требования версии Affero)
		\item пользователь имеет право знать, как обрабатываются его данные
	\end{itemize}
	
	А из этого непосредственно следует:
	
	\begin{itemize}
		\item любой может улучшить генератор
		\item любой может развернуть свою собственную копию для школы или университета
	\end{itemize}
	
	Текст и презентация защищены лицензией \textbf{CC-BY-SA 4.0} \cite{license_cc_by_sa}:
	
	\begin{itemize}
		\item Разрешено делиться, копировать и распространять презентацию и данную работу
		\item Можно изменять, трансформировать и брать за основу презентацию или данную работу для создания чего-то нового
		\item Лицензия не запрещает коммерческое использование
	\end{itemize}
	
	При этом необходимо указывать авторство и распространять все изменения под той же самой лицензией, что и оригинал.
	
	
	\newpage
	
	% --- СПИСОК ЛИТЕРАТУРЫ ---
	\addcontentsline{toc}{section}{Список литературы}
	\bibliographystyle{plainurl}
	\bibliography{references}
	
\end{document}